\section{Metode Validasi dan \textit{Benchmark} Klasik}
\label{sec:4_9_validasi_benchmark}

Algoritma VQE yang diimplementasikan pada penelitian ini merupakan sebuah \textit{proof of concept} yang memerlukan prosedur validasi ketat guna memastikan bahwa solusi portofolio yang dihasilkan selaras dengan pencarian solusi eksak. Prosedur validasi tersebut melibatkan metode \textit{brute force} yang mengevaluasi seluruh ruang keadaan kombinatorial secara deterministik. Selain itu, digunakan pula algoritma \textit{Sequential Least Squares Programming} (SLSQP) sebagai \textit{benchmark} performa industri untuk membandingkan efektivitas optimasi kontinu klasik terhadap pendekatan kuantum hibrida.

\subsection{Validasi Eksak melalui Algoritma \textit{Brute Force}}

Metode \textit{brute force} klasik bekerja dengan mengevaluasi nilai ekspektasi energi dari Hamiltonian total $\hat{\mathcal{H}}$ untuk setiap kemungkinan kombinasi portofolio biner dalam ruang Hilbert. Sebagai ilustrasi fisis, dilakukan evaluasi energi pada sampel data jendela Januari 2021 untuk sistem $N=2$ (BBCA dan TLKM) dengan parameter penalti $A = 4,1972$. Fungsi energi objektif $E(\vec{x})$ dihitung pada seluruh $2^N = 4$ keadaan basis sebagai berikut:
\begin{enumerate}
    \item \textbf{Keadaan $|00\rangle$ (Tanpa Aset):} 
    $E = \text{Obj}(0,0) + A(0+0-1)^2 = 0 + 4,1972(1) = 4,1972$.
    \item \textbf{Keadaan $|01\rangle$ (Hanya TLKM):} 
    $E = \text{Obj}(0,1) + A(0+1-1)^2 = -0,1183 + 0 = -0,1183$.
    \item \textbf{Keadaan $|10\rangle$ (Hanya BBCA):} 
    $E = \text{Obj}(1,0) + A(1+0-1)^2 = -0,1909 + 0 = -0,1909$.
    \item \textbf{Keadaan $|11\rangle$ (Kedua Aset):} 
    $E = \text{Obj}(1,1) + A(1+1-1)^2 \approx -0,29 + 4,1972 = 3,9072$.
\end{enumerate}
Hasil evaluasi tersebut menunjukkan bahwa keadaan $|10\rangle$ memiliki tingkat energi paling rendah, yang mana mengonfirmasi ketepatan sirkuit VQE dalam memusatkan amplitudo probabilitas pada konfigurasi aset yang paling optimal secara teoretis.

\subsection{Optimasi Kontinu melalui Algoritma SLSQP}

Algoritma \textit{Sequential Least Squares Programming} (SLSQP) merupakan metode optimasi numerik berbasis \textit{Sequential Quadratic Programming} (SQP) yang dirancang untuk menyelesaikan masalah minimisasi non-linear dengan batasan kesamaan (\textit{equality}) maupun ketidaksamaan (\textit{inequality}). Dalam konteks manajemen portofolio, SLSQP digunakan untuk menyelesaikan model Markowitz kontinu dengan formulasi fungsi biaya sebagai berikut:
\begin{equation}
    \min_{w} f(w) = \left[ \frac{\gamma}{2} w^T \Sigma w - \mu^T w \right]
\end{equation}
dengan batasan $\sum w_i = 1$ (anggaran penuh) dan $0 \leq w_i \leq 1$ (\textit{no short-selling}).

Secara teknis, SLSQP beroperasi dengan melakukan linierisasi pada batasan dan menerapkan aproksimasi kuadratik pada fungsi objektif melalui ekspansi deret Taylor orde kedua. Pada setiap iterasi $k$, algoritma menyusun sub-masalah \textit{Quadratic Programming} (QP) untuk menentukan arah pencarian optimal $d_k$ dengan bentuk:
\begin{equation}
    \min_{d} \nabla f(w_k)^T d + \frac{1}{2} d^T B_k d
\end{equation}
di mana $B_k$ merupakan aproksimasi matriks Hessian dari fungsi Lagrangian yang diperbarui menggunakan skema \textit{Broyden-Fletcher-Goldfarb-Shanno} (BFGS). Setelah arah $d_k$ diperoleh, algoritma melakukan pembaruan bobot $w_{k+1} = w_k + \alpha_k d_k$, dengan $\alpha_k$ sebagai ukuran langkah yang ditentukan melalui prosedur \textit{line search}.

\begin{algorithm}[H]
\caption{Optimasi Portofolio Klasik SLSQP}
\begin{algorithmic}[1]
\State \textbf{Input:} Matriks Kovariansi $\Sigma$, Imbal Hasil $\mu$, Parameter Risiko $\gamma$
\State \textbf{Inisialisasi:} Bobot awal $w_0$ (seragam), Matriks Hessian $B_0 = \mathbf{I}$
\While{tidak konvergen (\textit{tol} $< 10^{-6}$)}
    \State Hitung Gradien $\nabla f(w_k) = \gamma \Sigma w_k - \mu$
    \State Susun sub-masalah QP dengan batasan linear $A w = b$
    \State Selesaikan QP untuk mendapatkan arah pencarian $d_k$
    \State Tentukan ukuran langkah $\alpha_k$ melalui \textit{line search}
    \State Perbarui bobot: $w_{k+1} \leftarrow w_k + \alpha_k d_k$
    \State Perbarui aproksimasi Hessian $B_{k+1}$ menggunakan formula BFGS
\EndWhile
\State \textbf{Output:} Bobot optimal kontinu $w^*$
\State \textbf{Diskritisasi:} Pilih $K$ aset dengan $w^*$ tertinggi (\textit{Top-K Selection})
\end{algorithmic}
\end{algorithm}

Prosedur diskritisasi melalui \textit{Top-K Selection} diperlukan karena sirkuit VQE beroperasi pada domain biner, sementara SLSQP menghasilkan bobot dalam spektrum kontinu $[0,1]$. Dengan memilih $K$ aset yang memiliki alokasi bobot tertinggi, perbandingan antara metode klasik dan kuantum dapat dilakukan pada parameter batasan yang identik (\textit{fair comparison}). Hasil akhir dari algoritma ini berfungsi sebagai tolok ukur efisiensi alokasi modal yang akan dibandingkan dengan performa portofolio hasil optimasi VQE pada bab berikutnya.
